% GENERATED FILE. Edit canonical lessons, then run npm run generate:books.
% canonical-source-hash: fnv1a64:2596850f9648e1c4
% canonical-lessons: KA-C62-thread, KA-C62-needle, KA-C62-rope, KA-C62-broom, KA-C62-umbrella

\chapter{Things People Make}
\label{ch:ka-made}
\hlchaptermodality{62}

\begin{chapteropening}
\textbf{By the end of this chapter:} \emph{I can name five things a Kannada household keeps to hand, and ask politely for one.}

Complete KA-C62-umbrella and say all five words in order.
\end{chapteropening}

\section[nūlu]{\kn{ನೂಲು} (nūlu) — thread, yarn}

\label{lesson:KA-C62-thread}



\begin{quote}
Before the new run: what did \emph{bella} mean, and what did \emph{sakkare} mean?
\end{quote}

\subsection*{You'll want to know: \kn{ನೂಲು}}
\textbf{\kn{ನೂಲು}} (\emph{nūlu}) — \textquotedblleft{}thread, yarn\textquotedblright{}.

\kn{ನೂಲು} is inherited, and it is the stuff \kn{ಬಟ್ಟೆ} is made of. Everything worn in the cloth chapter began here.

The same shape does double duty: \kn{ನೂಲು} is the thread and it is equally the spinning of it, so one word covers the material and the making. Kannada does that often, and you will start to expect it.

\begin{grammarlens}[title={What you've built}]
The first of five things a household makes or keeps.
\end{grammarlens}

\subsection*{Your turn}
Say these aloud:

\begin{itemize}
\raggedright
  \item \emph{nūlu}
  \item it once more, slowly
  \item \emph{nūlu}, then \emph{baṭṭe}, and say which one comes first
\end{itemize}

\begin{tcolorbox}[breakable,colback=teal!4,colframe=teal!35!black,title={Before you move on}]
What does \kn{ನೂಲು} mean? (\textquotedblleft{}thread, yarn\textquotedblright{}.) Say it once more.
\end{tcolorbox}
\section[sūji]{\kn{ಸೂಜಿ} (sūji) — a needle}

\label{lesson:KA-C62-needle}



\begin{quote}
Before the new one: what did \emph{nūlu} mean?
\end{quote}

\subsection*{You'll want to know: \kn{ಸೂಜಿ}}
\textbf{\kn{ಸೂಜಿ}} (\emph{sūji}) — \textquotedblleft{}a needle\textquotedblright{}.

\kn{ಸೂಜಿ} is what the \kn{ನೂಲು} is put through. It is a borrowing from the north, and it arrived early enough that nothing about it feels foreign now.

It is also the hand of a clock. A \kn{ಗಂಟೆ} has two \kn{ಸೂಜಿ} on its face, and Kannada saw the resemblance before anybody thought to give the hands a separate name.

\begin{grammarlens}[title={What you've built}]
Two.
\end{grammarlens}

\subsection*{Your turn}
Say these aloud:

\begin{itemize}
\raggedright
  \item \emph{sūji}
  \item it once more, slowly
  \item \emph{nūlu}, then \emph{sūji}, and say which one goes through the other
\end{itemize}

\begin{tcolorbox}[breakable,colback=teal!4,colframe=teal!35!black,title={Before you move on}]
What does \kn{ಸೂಜಿ} mean? (\textquotedblleft{}a needle\textquotedblright{}.) Say it once more.
\end{tcolorbox}
\section[hagga]{\kn{ಹಗ್ಗ} (hagga) — a rope}

\label{lesson:KA-C62-rope}



\begin{quote}
Before the new one: what did \emph{sūji} mean?
\end{quote}

\subsection*{You'll want to know: \kn{ಹಗ್ಗ}}
\textbf{\kn{ಹಗ್ಗ}} (\emph{hagga}) — \textquotedblleft{}a rope\textquotedblright{}.

\kn{ಹಗ್ಗ} is \kn{ನೂಲು} grown up: strands laid together until the result will hold weight. Native, with nothing borrowed anywhere in it.

In a \kn{ಹಳ್ಳಿ} yard it is the thing at the end of which the \kn{ಹಸು} stands. Twist it thinner and it ties a \kn{ಬುಟ್ಟಿ} shut; twist it thicker and it draws water out of a \kn{ಕೆರೆ}.

\begin{grammarlens}[title={What you've built}]
Three.
\end{grammarlens}

\subsection*{Your turn}
Say these aloud:

\begin{itemize}
\raggedright
  \item \emph{hagga}
  \item it once more, slowly
  \item \emph{nūlu}, then \emph{hagga}, and say which one holds a cow
\end{itemize}

\begin{tcolorbox}[breakable,colback=teal!4,colframe=teal!35!black,title={Before you move on}]
What does \kn{ಹಗ್ಗ} mean? (\textquotedblleft{}a rope\textquotedblright{}.) Say it once more.
\end{tcolorbox}
\section[porake]{\kn{ಪೊರಕೆ} (porake) — a broom}

\label{lesson:KA-C62-broom}



\begin{quote}
Before the new one: what did \emph{hagga} mean?
\end{quote}

\subsection*{You'll want to know: \kn{ಪೊರಕೆ}}
\textbf{\kn{ಪೊರಕೆ}} (\emph{porake}) — \textquotedblleft{}a broom\textquotedblright{}.

\kn{ಪೊರಕೆ} is a bundle of stiff ribs tied at one end — no handle, so it is used bent over, in short strokes.

It is the first object of the day. The ground outside the \kn{ಬಾಗಿಲು} is swept with a \kn{ಪೊರಕೆ} before the \kn{ರಂಗೋಲಿ} is set out on it, because the pattern needs clean ground to sit on. Sweeping, then drawing: that order is fixed.

\begin{grammarlens}[title={What you've built}]
Four.
\end{grammarlens}

\subsection*{Your turn}
Say these aloud:

\begin{itemize}
\raggedright
  \item \emph{porake}
  \item it once more, slowly
  \item \emph{porake}, then \emph{raṅgōli}, and say which one comes first
\end{itemize}

\begin{tcolorbox}[breakable,colback=teal!4,colframe=teal!35!black,title={Before you move on}]
What does \kn{ಪೊರಕೆ} mean? (\textquotedblleft{}a broom\textquotedblright{}.) Say it once more.
\end{tcolorbox}
\section[koḍe]{\kn{ಕೊಡೆ} (koḍe) — an umbrella}

\label{lesson:KA-C62-umbrella}



\begin{quote}
Before the new one: what did \emph{porake} mean?
\end{quote}

\subsection*{You'll want to know: \kn{ಕೊಡೆ}}
\textbf{\kn{ಕೊಡೆ}} (\emph{koḍe}) — \textquotedblleft{}an umbrella\textquotedblright{}.

\kn{ಕೊಡೆ} works twice over in Karnataka: it keeps the \kn{ಮಳೆಗಾಲ} off you, and it keeps the \kn{ಬಿಸಿಲು} off you in the months on either side of it. In a Kannada street both uses are ordinary.

Read it carefully. \kn{ಕೊಡೆ} and the \kn{ಕೊಡು} you already know start with the same two letters, and they are separate items — one is a thing, the other is what you do with it, and joining them would be inventing a story. Kannada gives you plenty of real ones; this is not one.

Put \kn{ದಯವಿಟ್ಟು} in front of any of this chapter's five and you have asked for it properly.

\begin{grammarlens}[title={What you've built}]
Five: \kn{ನೂಲು}, \kn{ಸೂಜಿ}, \kn{ಹಗ್ಗ}, \kn{ಪೊರಕೆ}, \kn{ಕೊಡೆ}. Enough to ask a household politely for the thing you need.
\end{grammarlens}

\subsection*{Your turn}
Say these aloud:

\begin{itemize}
\raggedright
  \item \emph{koḍe}
  \item it once more, slowly
  \item all five in order, then \emph{dayaviṭṭu} in front of the one you want
\end{itemize}

\begin{tcolorbox}[breakable,colback=teal!4,colframe=teal!35!black,title={Before you move on}]
What does \kn{ಕೊಡೆ} mean? (\textquotedblleft{}an umbrella\textquotedblright{}.) Say it once more.
\end{tcolorbox}
